\chapter{Тормоза}
% полная версия: chapters/03_gaba.tex

Добавим в машину второй по численности тип нейронов --- \nt{ГАМК}. Их сигнал не подталкивает соседа к щелчку, а отталкивает. Это тормоза. Их немного, в коре от пятой части до четверти всех нейронов, но без них машина не ездит.

Есть одно ограничение: «плюсовой» нейрон не может сам затормозить соседа. Если нужна отрицательная связь, её приходится строить через посредника: «плюс» будит тормозной нейрон, а тот уже тормозит цель. Смена знака стоит одного нейрона и одной небольшой задержки.

\section*{«Но не если»}

Самая полезная операция с тормозами --- запрет: «сработай от A, но не если пришло B». Нейрон получает толчок от A и тормоз от посредника, которого будит B. Из запретов и «или» можно собрать любую логику.

\pencilfig{3}{\draw[pen=pcBlue] (0,0.6) .. controls (0.8,0.7) and (1.7,0.5) .. (2.55,0.6);
\draw[arr=pcBlue] (2.5,0.6) -- (2.62,0.6);
\draw[pen=pcBlue, wash=pcBlue] (3.0,0.6) circle (0.35);
\draw[pen=pcRed, wash=pcRed] (1.6,1.65) circle (0.28);
\draw[pen=pcGray] (1.6,2.45) -- (1.6,1.95); \draw[arr] (1.6,2.05) -- (1.6,1.95);
\draw[pen=pcRed] (1.78,1.43) -- (2.42,0.78);
\draw[pcRed, line width=1.3pt, line cap=round] (2.3,0.66) -- (2.56,0.92);
\node[lbl] at (-0.2,0.6) {А}; \node[lbl] at (1.6,2.65) {Б};}{«А, но не если Б»: красный нейрон перекрывает дорогу синему.}

Запрет во времени даёт детектор движения, который есть у вас в глазу. Два соседних датчика: один возбуждает выходной нейрон, другой тормозит его с задержкой. Если пятно света бежит в одну сторону с подходящей скоростью, торможение приходит ровно тогда, когда и возбуждение, --- и гасит его. Если в другую сторону --- торможение опаздывает, и нейрон успевает щёлкнуть. Получается нейрон, который видит движение только в одну сторону.

\section*{Вычитание или деление}

Торможение умеет не только вычитать. Открытый тормозной канал не тянет напряжение вниз, он тянет его \emph{к покою} --- как будто в ведро проделали ещё одну дырку. От этого любой приток воды поднимает её меньше. Торможение не отнимает от сигнала постоянную порцию, а \emph{делит} его: это регулятор громкости. Если сила торможения растёт вместе с общей активностью вокруг, каждый нейрон реагирует не на то, насколько силён его сигнал, а на то, насколько он силён \emph{по сравнению с соседями}. Поэтому одна и та же сеть работает и в полумраке, и на ярком солнце. Правда, на частоту импульсов такой «делитель» действует чисто только вместе с постоянным фоновым шумом, который в живом мозге есть всегда.

Есть и второе следствие: лишняя дырка заставляет ведро быстрее забывать. Окно, в которое два сигнала должны попасть, чтобы сложиться, сужается. Торможение делает детекторы совпадений строже.

\section*{Выбор}

Теперь главное, чего не хватало: выбор одного из многих. Две группы нейронов тормозят друг друга. Пока торможение слабое, работают обе, просто тише. Но если оно достаточно сильное, любая маленькая разница нарастает, как на качелях, пока одна группа не замолчит совсем. Победитель забирает всё. И победа держится: чтобы сеть передумала, новый довод должен быть заметно сильнее старого. Машина не дёргается от каждого шороха.

Тормоза стирают и память: тормозной нейрон по сигналу «сброс» гасит костёр из прошлой главы. Две такие ячейки, тормозящие друг друга, --- аналог защёлки, основы любой компьютерной памяти.

\section*{Жизнь на грани}

Сеть из одних «плюсов» может уйти в лавину: каждый импульс рождает больше одного следующего. Тормоза держат её в рабочей точке, как регулятор держит температуру в духовке. Есть два условия: торможение должно быть достаточно сильным и достаточно быстрым. Если оно сильное, но медленное, регулятор начинает перебарщивать --- как человек, который крутит кран в душе, не дождавшись, пока дойдёт горячая вода. Температура качается туда-сюда. В мозге такая петля из «плюсов» и «минусов» порождает быстрый ритм --- несколько десятков колебаний в секунду.

Кора, по-видимому, живёт именно на грани: её собственная обратная связь достаточно сильна, чтобы без тормозов уйти в лавину. Удерживают её только быстрые тормоза, и поэтому она так чутко отвечает на слабые сигналы.

Вывод главы неожиданный: тормоза почти не добавляют машине новых вычислений. Они добавляют \emph{управление}: выбор, сброс, регулировку громкости, устойчивость. Возбуждение несёт данные, торможение решает, каким данным пройти, когда и насколько громко.

\fullversion{3}
